\chapter{The One-Particle-Irreducible Renormalization Group}
\label{ch:volume-ii-11}

\section{The 1PI Effective Action}

The one-particle-irreducible effective action \(\Gamma[\varphi]\) is the
Legendre transform of the connected generating functional. Its functional
derivatives generate 1PI vertices:
\[
  \Gamma^{(n)}(x_1,\ldots,x_n)
  =
  \frac{\delta^n\Gamma}{\delta\varphi(x_1)\cdots\delta\varphi(x_n)}.
\]
In perturbation theory, \(\Gamma\) is the sum of 1PI diagrams evaluated in a
background field. It packages the inverse exact propagator, proper vertices,
and local counterterms in one object.

\section{Renormalization Conditions}

A momentum-subtraction scheme fixes couplings by specifying values of 1PI
vertices at a scale \(\mu\). For example, in a scalar theory,
\[
  \Gamma^{(2)}(p)\big|_{p^2=\mu^2},\qquad
  \Gamma^{(4)}(p_i)\big|_{\mathcal S_\mu}
\]
may define the renormalized mass, field normalization, and quartic coupling.
Changing \(\mu\) while holding bare quantities fixed produces the 1PI RG
equations.

\section{Improving Perturbation Theory}

The 1PI RG is especially useful when logarithms become large. A fixed-order
amplitude may contain terms such as
\[
  g^2\log\frac{p}{\mu}.
\]
Choosing \(\mu\sim p\), or equivalently running the coupling to the physical
scale, resums families of logarithms. The RG does not create new physics from
nothing; it reorganizes perturbation theory around the correct scale.

\section{Limits Of The 1PI View}

The 1PI effective action is not the Wilsonian effective action. It includes
effects of all momentum scales and is generally nonlocal as a functional. Its
Euclidean version is convex when defined exactly, a fact that can obscure
spontaneous symmetry breaking if interpreted naively. The Wilsonian action,
introduced later, is scale-local and better suited to questions about
integrating out modes.
